\begin{name}
	{\tenchude}
	{ÔN TẬP CHƯƠNG V}
	{LỚP TOÁN THẦY PHÁT}
	{Thời gian: 90 phút - Không kể thời gian phát đề}
\end{name}
\Opensolutionfile{ans}[ans/ans-2-OTC5-TN]
\caulc
%%%============EX_1==============%%%
\begin{ex}%[2H5N3-2]
	Trong không gian $Oxyz$, cho mặt cầu $(S)\colon (x+3)^2+(y+1)^2+(z-1)^2=2$. Xác định tọa độ tâm của mặt cầu $(S)$.
	\choice
	{$I(-3; 1;-1)$}
	{$I(3; 1;-1)$}
	{\True $I(-3;-1; 1)$}
	{$I(3;-1; 1)$}
	\loigiai{Mặt cầu $(S)$ có tâm là $I(-3;-1; 1)$}
\end{ex}
%%%==============EX_2============%%%
\begin{ex}%[2H5H3-2]
	Trong không gian với hệ tọa độ $Oxyz$, cho mặt cầu $(S)\colon x^2+y^2+z^2-8x+2y+1=0$. Tìm tọa độ tâm và bán kính của mặt cầu $(S)$.
	\choice
	{$I(-4; 1; 0), R=2$}
	{$I(-4; 1; 0), R=4$}
	{$I(4;-1; 0), R=2$}
	{\True $I(4;-1; 0), R=4$}
	\loigiai{
		Ta có  $x^2+y^2+z^2-8x+2y+1=0\Leftrightarrow(x-4)^2+(y+1)^2+z^2=16$.\\
		Vậy mặt cầu $(S)$ có tâm $I(4;-1; 0)$ và bán kính $R=4$.}
\end{ex}
%%%==============EX_3============%%%
\begin{ex}%[2H5H3-3]
	Trong không gian với hệ tọa độ $Oxyz$, cho hai điểm $A(1;-2; 7)$, $B(-3; 8;-1)$. Mặt cầu đường kính $AB$ có phương trình là
	\choice
	{$(x+1)^2+(y-3)^2+(z-3)^2=\sqrt{45}$}
	{$(x-1)^2+(y+3)^2+(z+3)^2=45$}
	{$(x-1)^2+(y-3)^2+(z+3)^2=\sqrt{45}$}
	{\True $(x+1)^2+(y-3)^2+(z-3)^2=45$}
	\loigiai{
		Gọi $I$ là trung điểm $AB$ ta có $I(-1; 3; 3)$ là tâm mặt cầu.\\
		Bán kính $R=IA=\sqrt{(1+1)^2+(-2-3)^2+(7-3)^2}=\sqrt{45}$.\\
		Vậy phương trình mặt cầu cần tìm là $(x+1)^2+(y-3)^2+(z-3)^2=45$.
	}
\end{ex}
%%%==============EX_4============%%%
\begin{ex}%[2H5N2-2]
	Vectơ nào sau đây là một vectơ chỉ phương của đường thẳng $\dfrac{x+2}{3}=\dfrac{y+1}{-2}=\dfrac{z-3}{-1}$ 
	\choice
	{$(-2; 1;-3)$}
	{\True $(-3; 2; 1)$}
	{$(3;-2; 1)$}
	{$(2; 1; 3)$}
	\loigiai{Vectơ chỉ phương của đường thẳng là $\vec{u}=(3;-2;-1)=-1(-3; 2; 1)$ nên $\vec{u_1}=(-3; 2; 1)$ cũng là một vectơ chỉ phương của đường thẳng.
	}
\end{ex}
%%%==============EX_5============%%%
\begin{ex}%[2H5H2-3]
	Trong không gian với hệ tọa độ $Oxyz$, cho hai điểm $M(1;-2; 1)$, $N(0; 1; 3)$. Phương trình đường thẳng qua hai điểm $M$, $N$ là
	\choice
	{$\dfrac{x+1}{-1}=\dfrac{y-2}{3}=\dfrac{z+1}{2}$}
	{$\dfrac{x+1}{1}=\dfrac{y-3}{-2}=\dfrac{z-2}{1}$}
	{\True $\dfrac{x}{-1}=\dfrac{y-1}{3}=\dfrac{z-3}{2}$}
	{$\dfrac{x}{1}=\dfrac{y-1}{-2}=\dfrac{z-3}{1}$}
	\loigiai{
		Ta có  $\vec{MN}=(-1; 3; 2)$.\\
		Đường thẳng $MN$ qua $N(0; 1; 3)$ nhận $\vec{MN}=(-1; 3; 2)$ làm vectơ chỉ phương có phương trình $$\dfrac{x}{-1}=\dfrac{y-1}{3}=\dfrac{z-3}{2}.$$
	}
\end{ex}
%%%==============EX_6============%%%
\begin{ex}%[2H5H2-5]
	Trong không gian với hệ tọa độ $Oxyz$, cho đường thẳng $d\colon \heva{&x=1+2t \\ &y=3-t \\ &z=1-t}$, $t \in \mathbb{R}$ và mặt phẳng $(P)\colon x+2y-3z+2=0$. Tìm tọa độ của điểm $A$ là giao điểm của đường thẳng $d$ và mặt phẳng $(P)$.
	\choice
	{$A(3; 5; 3)$}
	{$A(1; 3; 1)$}
	{\True $A(-3; 5; 3)$}
	{$A(1; 2;-3)$}
	\loigiai{
		Vì $A$ là giao điểm của đường thẳng $d$ và mặt phẳng $(P)$ nên $A\in d \Rightarrow A(1+2t; 3-t; 1-t)$.\\
		Mà $A\in(P) \Rightarrow(1+2t)+2(3-t)-3(1-t)+2=0\Rightarrow t=-2$.\\
		Vậy tọa độ điểm $A(-3; 5; 3)$.
	}
\end{ex}
%%%==============EX_7============%%%
\begin{ex}%[2H5H2-7]
	Trong không gian với hệ tọa độ $Oxyz$ cho đường thẳng $d\colon \heva{&x=1-t \\&y=2+2t \\&z=3+t
	}$ và mặt phẳng $(P)\colon x-y+3=0$. Tính số đo góc giữa đường thẳng $d$ và mặt phẳng $(P)$.
	\choice
	{\True $60^{\circ}$}
	{$30^{\circ}$}
	{$120^{\circ}$}
	{$45^\circ$}
	\loigiai{
		Đường thẳng $d$ có vectơ chỉ phương là $\vec{u}=(-1; 2; 1)$.\\
		Mặt phẳng $(P)$ có vectơ pháp tuyến là $\vec{n}=(1;-1; 0)$.\\
		Gọi $\alpha$ là góc giữa đường thẳng $d$ và mặt phẳng $(P)$. \\
		Khi đó ta có
		$$\sin \alpha=\dfrac{|\vec{u} \cdot \vec{n}|}{|\vec{u}|\cdot |\vec{n}|}=\dfrac{|-1 \cdot 1+2 \cdot(-1)+1 \cdot 0|}{\sqrt{(-1)^2+2^2+1^2} \cdot \sqrt{1^2+(-1)^2+0^2}}=\dfrac{3}{2 \sqrt{3}}=\dfrac{\sqrt{3}}{2}.$$
		Do đó $\alpha=60^{\circ}$.
	}
\end{ex}
%%%==============EX_8============%%%
\begin{ex}%[2H5H2-6]
	Trong không gian $Oxyz$, cho hai đường thẳng $d_1\colon  \dfrac{x-1}{2}=\dfrac{y+2}{-1}=\dfrac{z}{1}$ và $d_2\colon \heva{&x=1+4t \\&y=-1-2t \\&z=2+2t&}$.
	Khoảng cách giữa hai đường thẳng đã cho bằng
	\choice
	{$\dfrac{\sqrt{87}}{6}$}
	{\True $\dfrac{\sqrt{174}}{6}$}
	{$\dfrac{\sqrt{174}}{3}$}
	{$\dfrac{\sqrt{87}}{3}$}
	\loigiai{
		Đường thẳng $d_1$ đi qua điểm $M(1;-2; 0)$ và nhận $\vec{u_1}=(2;-1; 1)$ làm VTCP.\\
		Đường thẳng $d_2$ đi qua điểm $N(1;-1; 2)$ và nhận $\vec{u_2}=(4;-2; 2)$ làm VTCP.\\
		Vì $\vec{u_2}=2\cdot \vec{u_1}$ nên đường thẳng $d_1$ song song hoặc trùng với đường thẳng $d_2$.\\
		Lại có điểm $M(1;-2; 0) \in d_1$ nhưng $M(1;-2; 0) \notin d_2$ nên suy ra $d_1 \parallel d_2$.\\
		Vậy khoảng cách giữa hai đường thẳng đã cho bằng khoảng cách từ điểm $M(1;-2; 0)$ đến đường thẳng $d_2$.\\
		$$\mathrm{d}\left(M; d_2\right)=\dfrac{\left| \left[\vec{MN}; \vec{u_2}\right]\right| }{\left|\vec{u_2}\right|}.$$
		Ta có $\vec{MN}=(0; 1; 2),\left[\vec{MN}; \vec{u_2}\right]=(6; 8;-4)$.\\
		$\Rightarrow d\left(d_1; d_2\right)=d\left(M; d_2\right)=\dfrac{\sqrt{6^2+8^2+(-4)^2}}{\sqrt{4^2+(-2)^2+2^2}}=\dfrac{\sqrt{174}}{6}$.
	}
\end{ex}
%%%==============EX_9============%%%
\begin{ex}%[2H5H2-5]
	Trong không gian $Oxyz$, cho hai đường thẳng chéo nhau ${d_1\colon  \dfrac{x-2}{2}=\dfrac{y-6}{-2}=\dfrac{z+2}{1}}$ và ${d_2\colon  \dfrac{x-4}{1}=\dfrac{y+1}{3}=\dfrac{z+2}{-2}}$.Phương trình mặt phẳng $(P)$ chứa $d_1$ và $(P)$ song song với đường thẳng $d_2$ là
	\choice
	{\True $(P)\colon x+5y+8z-16=0$}
	{$(P)\colon x+5y+8z+16=0$}
	{$(P)\colon x+4y+6z-12=0$}
	{$(P)\colon 2x+y-6=0$}
	\loigiai{
		Đường thẳng $d_1$ đi qua $A(2; 6;-2)$ và có một vectơ chỉ phương $\vec{u_1}=(2;-2; 1)$.\\
		Đường thẳng $d_2$ có một vectơ chỉ phương $\vec{u_2}=(1; 3;-2)$.\\
		Gọi $\vec{n}$ là một vectơ pháp tuyến của mặt phẳng $(P)$. Do mặt phẳng $(P)$ chứa $d_1$ và $(P)$ song song với đường thẳng $d_2$ nên $\vec{n}=\left[\vec{u_1}, \vec{u_2}\right]=(1; 5; 8)$.\\
		Vậy phương trình mặt phẳng $(P)$ đi qua $A(2; 6;-2)$ và có một vectơ pháp tuyến $\vec{n}=(1; 5; 8)$ là $x+5y+8z-16=0$.
	}
\end{ex}
%%%==============EX_10============%%%
\begin{ex}%[2H5N1-3]
	Trong không gian $Oxyz$, mặt phẳng $(P)$ đi qua điểm $M(3;-1; 4)$ đồng thời vuông góc với giá của vectơ $\vec{a}=(1;-1; 2)$ có phương trình là
	\choice
	{$3x-y+4z-12=0$}
	{$3x-y+4z+12=0$}
	{\True $x-y+2z-12=0$}
	{$x-y+2z+12=0$}
	\loigiai{
		$(P)$ có phương trình là $1\cdot(x-3)-1\cdot(y+1)+2\cdot(z-4)=0\Leftrightarrow x-y+2z-12=0$.}
\end{ex}
%%%==============EX_11============%%%
\begin{ex}%[2H5N1-5]
	Trong không gian $Oxyz$, tính khoảng cách từ $M(1; 2;-3)$ đến mặt phẳng $(P)\colon x+2y+2z-10=0$.
	\choice
	{\True $\dfrac{11}{3}$}
	{$3$}
	{$\dfrac{7}{3}$}
	{$\dfrac{4}{3}$}
	\loigiai{
		$$\mathrm{d}(M;(P))=\dfrac{|1+2\cdot 2+2\cdot(-3)-10|}{\sqrt{1^2+2^2+2^2}}=\dfrac{|-11|}{3}=\frac{11}{3}.$$}
\end{ex}
%%%==============EX_12============%%%
\begin{ex}%[2H5H1-3]
	Trong không gian với hệ tọa độ $Oxyz$, cho ba điểm $A(2; 0; 0), B(0;-1; 0), C(0; 0;-3)$. Viết phương trình mặt phẳng $(ABC)$.
	\choice
	{$-3x+6y-2z+6=0$}
	{$-3x-6y+2z+6=0$}
	{\True $-3x+6y+2z+6=0$}
	{$-3x-6y+2z-6=0$}
	\loigiai{
		Phương trình mặt phẳng $(ABC)$ là $\dfrac{x}{2}+\dfrac{y}{-1}+\dfrac{z}{-3}=1\Leftrightarrow-3x+6y+2z+6=0$.}
\end{ex}
\Closesolutionfile{ans}
\indapan{6}{ans/ans-2-OTC5-TN}

\cauds
\Opensolutionfile{ans}[ans/ans-2-OTC5-DS]
%%%==============EX_1============%%%
\begin{ex}%[2H5V1-3]
	Trong không gian với hệ tọa độ $Oxyz$, cho ba điểm  $A(1;-2; 1)$, $B(-1; 3; 3)$, $C(2;-4; 2)$.
	\choiceTF
	{\True $\vec{AB}$ và $\vec{AC}$ là cặp vectơ chỉ phương của mặt phẳng $(ABC)$}
	{$\vec{n}=(9; 4; 1)$ là một vectơ pháp tuyến của mặt phẳng $(ABC)$}
	{\True $\vec{m}=(-9;-4; 1)$ là một vectơ pháp tuyến của mặt phẳng $(ABC)$}
	{\True Phương trình mặt phẳng $(ABC)$ là $9x+4y-z=0$}
	\loigiai{
		\begin{itemchoice}
			\itemch Đúng. Vì $A, B, C\in(ABC)$ nên $\vec{AB}$ và $\vec{AC}$ là cặp vectơ chỉ phương của mặt phẳng $(ABC)$.
			\itemch Sai. Ta có $\vec{AB}=(-2; 5; 2), \vec{AC}=(1;-2; 1) \Rightarrow \vec{n}=\left[\vec{AB}, \vec{AC}\right]=(9; 4;-1)$.
			\itemch Đúng. Ta có $\vec{n}=\left[\vec{AB}, \vec{AC}\right]=(9; 4;-1) \Rightarrow \vec{m}=-1\vec{n}=(-9;-4; 1)$ là một vectơ pháp tuyến của mặt phẳng $(ABC)$.
			\itemch Đúng. Ta có $(ABC)$ qua $A(1;-2; 1)$ và có vectơ pháp tuyến $\vec{n}=\left[\vec{AB}, \vec{AC}\right]=(9; 4;-1)\Rightarrow(ABC)\colon 9x+4y-z=0$.
		\end{itemchoice}
		
	}
\end{ex}
\begin{ex}%[2H5V2-4]
	Trong không gian với hệ tọa độ $Oxyz$, cho $\Delta\colon \heva{&x=t \\&y=1-2t \\&z=2},\,t\in \mathbb{R}$
	\choiceTF
	{\True Hệ phương trình trên được gọi là phương trình tham số của đường thẳng $\Delta$}
	{Điểm $A=(1;-2; 0)$ là một điểm thuộc đường thẳng $\Delta$}
	{\True $\vec{a}=(-2; 4; 0)$ là một vectơ chỉ phương của đường thẳng $\Delta$}
	{Đường thẳng $\Delta$ vuông góc với đường thẳng $d\colon \heva{&x=2+2t \\
			&y=1+2t \\
			&z=1-3t}$}
	\loigiai{
		\begin{itemchoice}
			\itemch Đúng. Vì đây là định nghĩa phương trình tham số của đường thẳng.
			\itemch Sai. Vì thay điểm $A$ vào phương trình tham số, ta nhận được 3 giá trị $t$ khác nhau.
			\itemch Đúng. Vì $\vec{u}=(1;-2; 0)$ là một vectơ chỉ phương của đường thẳng $\Delta$ nên $\vec{a}=-2\vec{u}=(-2; 4; 0)$ cũng là một vectơ chỉ phương của đường thẳng $\Delta$.
			\itemch Sai. Vì $1\cdot 2-2\cdot 2-0\cdot 3=-2\neq 0$.
		\end{itemchoice}
	}
\end{ex}
\begin{ex}%[2H5V2-7]
	Trong không gian với hệ tọa độ $Oxyz$, cho hai đường thẳng $\Delta_1\colon  \dfrac{x-1}{1}=\dfrac{y-1}{-1}=\dfrac{z-2}{-4}$ và $\Delta_2\colon \heva{&x=2t\\&y=1-2t \\&z=-1-8t.}$
	\choiceTF
	{\True Phương trình đường thẳng $\Delta_1$ được gọi là phương trình chính tắc}
	{\True$\vec{u}=(1;-1;-4)$ là một vectơ chỉ phương của đường thẳng $\Delta_1$}
	{$\Delta_1 \equiv \Delta_2$}
	{Góc giữa hai đường thẳng $\Delta_1$ và $\Delta_2$ bằng $30^{\circ}$}
	\loigiai{
		\begin{itemchoice}
			\itemch Đúng. Vì đây là định nghĩa phương trình chính tắc của đường thẳng.
			\itemch Đúng. Vì $\vec{u}=(1;-1;-4)$ là một vectơ chỉ phương của đường thẳng $\Delta_1$.
			\itemch Sai. Vì $M(1;1;2) \notin \Delta_2$ và $\vec{u_2}=2\vec{u_1}$ nên $\Delta_1 \parallel \Delta_2$.
			\itemch Sai. Vì $\Delta_1 \parallel \Delta_2$ nên góc giữa hai đường thẳng $\Delta_1$ và $\Delta_2$ bằng $0^{\circ}$.
		\end{itemchoice}
	}
\end{ex}
%%%==============EX_3============%%%
\begin{ex}%[2H5V3-3]
	Cho mặt cầu $(S)\colon  x^2+y^2+z^2-4x-4y-4z=0$ và điểm $A(1;-4; 3)$.
	\choiceTF
	{Mặt cầu $(S)$ có tâm $(-2;-2;-2)$}
	{\True Mặt cầu $(S)$ có bán kính $R=2\sqrt{3}$}
	{Điểm A nằm trong mặt cầu $(S)$}
	{\True Phương trình của mặt cầu $\left(S_1\right)$ có tâm $A$ và bán kính $R=3$ là $(x-1)^2+(y+4)^2+(z-3)^2=9$}
	\loigiai{
		\begin{itemchoice}
			\itemch Sai. Vì $(S)$ có tâm $I(2; 2; 2)$. 
			\itemch Đúng. Vì bán kính $R=\sqrt{2^2+2^2+2^2-0}=2\sqrt{3}$.
			\itemch Sai. Vì $IA=\sqrt{(1-2)^2+(-4-2)^2+(3-2)^2}=\sqrt{38} > 2\sqrt{3}$.
			\itemch Đúng. Vì Phương trình của mặt cầu $\left(S_1\right)$ có tâm $A$ và bán kính $R=3$ là
			$(x-1)^2+(y+4)^2+(z-3)^2=9$. 
		\end{itemchoice}
	}
\end{ex}
\Closesolutionfile{ans}
\indapan{4}{ans/ans-2-OTC5-DS}
\Opensolutionfile{ans}[ans/ans-2-OTC5-KQ]
\caukq
%%%==============EX_2============%%%
\begin{ex}%[2H5H3-3]
	Trong không gian với hệ trục tọa độ $Oxyz$, cho ba điểm $A(1; 0; 0), C(0; 0; 3), B(0; 2; 0)$. Tập hợp các điểm $M$ thỏa mãn $MA^2=MB^2+MC^2$ là mặt cầu có bán kính là bao nhiêu?
	\shortans{$1{,}41$}
	\loigiai{
		Giả sử $M(x; y; z)$.\\
		Ta có  $MA^2=(x-1)^2+y^2+z^2; \quad MB^2=x^2+(y-2)^2+z^2; \quad MC^2=x^2+y^2+(z-3)^2$.\\
		Khi đó 
		\begin{eqnarray*}
			MA^2=MB^2+MC^2&\Leftrightarrow&(x-1)^2+y^2+z^2=x^2+(y-2)^2+z^2+x^2+y^2+(z-3)^2\\
			&\Leftrightarrow&-2x+1=(y-2)^2+x^2+(z-3)^2 \\
			&\Leftrightarrow& (x+1)^2+(y-2)^2+(z-3)^2=2
		\end{eqnarray*}
		Vậy tập hợp các điểm $M$ thỏa mãn $MA^2=MB^2+MC^2$ là mặt cầu có bán kính là $R=\sqrt{2}$.
	}
\end{ex}
%%%==============EX_3============%%%
\begin{ex}%[2H5V3-3]
	Trong không gian $Oxyz$, gọi $I(a; b; c)$ là tâm mặt cầu đi qua điểm $A(1;-1; 4)$ và tiếp xúc với tất cả các mặt phẳng tọa độ $Oxy$, $Oyz$, $Ozx$. Tính $P=a-b+c$.
	\shortans{$9$}
	\loigiai{
		Vì mặt cầu tâm $I$ tiếp xúc với các mặt phẳng tọa độ nên
		$$d(I,(O y z))=d(I,(O z x))=d(I,(Oxy))=R \Leftrightarrow|a|=|b|=|c|=R$$
		Vì $A(1;-1; 4)$ nên $a>0$, $b<0$, $c>0$; do đó $I(R;-R; R)$.\\
		Mặt cầu tâm $I(R;-R; R)$ bán kính $R$ đi qua $A(1;-1; 4)$ nên
		$$(R-1)^2+(-R-1)^2+(R-4)^2=R^2\Leftrightarrow R^2-6R+9=0\Leftrightarrow R=3.$$
		Khi đó $P=a-b+c=9$.
	}
\end{ex}
%%%==============EX_4============%%%
\begin{ex}%[2H5V2-5]
	Trong không gian với hệ tọa độ $Oxyz$, cho điểm $M(-4; 0; 0)$ và đường thẳng $\Delta\colon \heva{&x=1-t \\
		&y=-2+3t \\&z=-2t}$. Gọi $H(a; b; c)$ là hình chiếu của $M$ lên $\Delta$. Tính $a+b+c$.
	\shortans{$-1$}
	\loigiai{
		$H$ là hình chiếu của $M$ lên $\Delta$ nên tọa độ của $H$ là nghiệm hệ phương trình
		$$\heva{&x=1-t \\&y=-2+3t \\&z=-2t \\& -x+3y-2z=-(-4)} \Leftrightarrow \heva{&x=\dfrac{3}{14}\\&y= \dfrac{5}{14}\\ &z= \dfrac{-22}{14}}$$
		Vậy $H\left(\dfrac{3}{14}; \dfrac{5}{14}; \dfrac{-22}{14}\right)$ nên $a+b+c=-1$.
	}
\end{ex}
%%%==============EX_5============%%%
\begin{ex}%[2H5V2-4]
	Trong không gian $Oxyz$, cho hình thang cân $ABCD$ có đáy là $AB$ và $CD$. Biết $A(3; 1;-2)$, $B(-1; 3; 2), C(-6; 3; 6)$ và $D (a; b; c)$ với $a, b, c \in \mathbb{R}$. Giá trị của $a+b+c$.
	\shortans{$-3$}
	\loigiai{
		\textbf{Cách 1:} Phương trình đường thẳng $d$ qua $C(-6; 3; 6)$ và song song với đường thẳng $AB$ là
		$$\dfrac{x+6}{-2}=\dfrac{y-3}{1}=\dfrac{z-6}{2}.$$
		Điểm $D$ thuộc đường thẳng $d$ nên gọi tọa độ $D$ là $D(-6-2t; 3+t; 6+2t)$.\\
		Tứ giác $ABCD$ là hình thang cân nên ta có
		$$|\vec{A D}|=|\vec{B C}|\Leftrightarrow t^2+8 t+12=0 \Leftrightarrow \hoac{& t=-2 \\ &t=-6}$$
		Với $t=-2\Rightarrow D_1(-2; 1; 2)$, tứ giác là hình bình hành nên loại.\\
		Với $t=-6\Rightarrow D_2(6;-3;-6)$ thỏa mãn, nên $6-3-6=-3$.\\
		\textbf{Cách 2:} Điểm $D$ thuộc đường thẳng $d$ qua $C(-6; 3; 6)$ và song song với đường thẳng $AB$ là
		$$\dfrac{x+6}{-2}=\dfrac{y-3}{1}=\dfrac{z-6}{2}.$$
		Điểm $H$ là hình chiếu vuông góc của $A$ lên $CD$ nên tọa độ $H$ là nghiệm hệ phương trình
		$$\heva{&x=-6-2t \\&y=3+t \\&z=6+2t \\& -2x+1y+2z=-2 \cdot 3+1 \cdot 1+2 \cdot (-2)} \Leftrightarrow \heva{&x=2\\&y=-1\\&z=-2}$$
		Tương tự điểm $K$ là hình chiếu vuông góc của $B$ lên $CD$ nên $K(-2; 1; 2)$.\\
		Vì $ABCD$ là hình thang cân nên $\vec{HD}=\vec{CK}\Rightarrow D(6;-3;-6)$ nên $a+b+c=-3$.
	}
\end{ex}
%%%==============EX_6============%%%
\begin{ex}%[2H5V3-4]
	Phía trước của một sân vận động, người ta có đặt một quả bóng đá khổng lồ. Cách đó không xa, có một bóng đèn. Trong không gian $Oxyz$, giả sử mặt ngoài quả bóng $(S)$ có phương trình $(x-1)^2+(y-1)^2+(z-3)^2=9$ và bóng đèn ở vị trí $M(15 ; 1 ; 6)$. Khi đó, khoảng cách xa nhất mà bóng đèn có thể chiếu đến một điểm trên quả bóng bằng $M A$, với $A$ là điểm nằm trên mặt cầu $(S)$. Vậy $M A$ bằng bao nhiêu mét?	
	\shortans{$14$}
	\loigiai{
		\begin{center}
			\begin{tikzpicture}[>=stealth,line width=1pt,scale=0.8]
				\def\r{3}
				\path
				(0,0) coordinate (I)
				(-8,0) coordinate (M)
				(I)+(115: \r) coordinate (A)
				(I)+(-115: \r) coordinate (B)
				;
				\draw (I) circle (\r);
				\draw (A)--(I)--(B)--(M)--cycle
				pic[draw, angle radius=2mm]{right angle=I--A--M}
				pic[draw, angle radius=2mm]{right angle=I--B--M}
				;
				\foreach \x/\g in {I/0,A/90,B/-90,M/-180}
				\draw [fill=white] (\x) circle (1pt) + (\g: 0.3) node {$\x$};
			\end{tikzpicture}	
		\end{center}
		Bề mặt ngoài của quả bóng $(S)$ là một mặt cầu có tâm $I(1 ; 1 ; 3)$ và bán kính $R=3$. \\
		Ta có $MI=\sqrt{(1-15)^2+(1-1)^2+(3-6)^2}=\sqrt{205}$.\\
		Suy ra $MA=\sqrt{MI^2-IA^2}=\sqrt{205-9}=\sqrt{196}=14$.\\
		Vậy $M A=14 \,\mathrm{m}$.		
	}
\end{ex}
\Closesolutionfile{ans}
\indapan{6}{ans/ans-2-OTC5-KQ}

